% SPDX-License-Identifier: GPL-3.0-or-later
% Copyright (C) 2026 Alessandro Greco
\chapter{Правовые и этические аспекты в Италии и ЕС}\label{cap:25-legale}
\aggiornatoal{23 сентября 2026 г.}{Регламент (ЕС) 2022/2065; Законодательный декрет 196/2003 (сводный текст, опубликованный Итальянским органом по защите персональных данных); рекомендации Tor Research Safety Board}

\begin{obiettivi}
К концу этой главы вы сможете:
\begin{itemize}[nosep]
  \item отличать законность инструмента от законности способа его использования;
  \item интерпретировать европейское правило \enquote{простая передача} (\emph{mere conduit}) применительно к работе выходного ретранслятора Tor;
  \item определять нерешённые вопросы для операторов ретрансляторов в Италии, включая ответственность, хранение данных и изъятие оборудования;
  \item применять принципы безопасности Tor Research Safety Board к исследованиям Tor.
\end{itemize}
\end{obiettivi}

\begin{attenzione}[Эта глава не является юридической консультацией]
Право меняется, зависит от обстоятельств конкретного дела, а решений, непосредственно касающихся Tor, относительно немного. Эта глава очерчивает соответствующие вопросы и указывает нормы, рассмотренные по их официальному тексту, прямо обозначая неопределённость. Любой, кто принимает эксплуатационное решение, например запускает выходной ретранслятор, должен получить консультацию юриста, специализирующегося на технологическом праве.
\end{attenzione}

% ---------------------------------------------------------------------------
\section{Инструмент и способы его использования}\label{sec:legale-strumento}

В Италии, как и в других странах Европейского союза, нет закона, который в общем случае запрещал бы использование Tor или других инструментов, защищающих конфиденциальность коммуникаций. Напротив, тайна переписки и иных сообщений (статья~15 Конституции Италии) и свобода выражения мнения (статья~21) являются основными правами, как и уважение частной жизни и защита персональных данных по статьям~7 и~8 Хартии Европейского союза об основных правах.

Действия, совершаемые через Tor, подчиняются тем же законам, что действуют и без Tor. Несанкционированный доступ к компьютерной системе (статья~615-\emph{ter} Уголовного кодекса Италии) не становится законным только потому, что проходит через три ретранслятора. Различие между инструментом и его использованием является центральным для этой главы и важно также для тех, кто предоставляет инфраструктуру: операторы ретрансляторов предоставляют возможность передачи данных, подобно провайдерам связи, а не выбирают передаваемое содержимое.

% ---------------------------------------------------------------------------
\section{Операторы ретрансляторов}\label{sec:legale-relay}

\subsection{Простая передача (mere conduit)}

Европейский \emph{Регламент о цифровых услугах} (Digital Services Act, DSA), применяемый с 17 февраля 2024 года, заменил правила ответственности посредников, первоначально установленные Директивой об электронной коммерции 2000 года, которую Италия реализовала Законодательным декретом 70/2003~\cite{dsa2022}. Статья~4 устанавливает режим \textbf{простой передачи} (\emph{mere conduit}): сервис, передающий информацию, предоставленную другими лицами, освобождается от ответственности за эту информацию, если он не инициирует передачу, не выбирает получателя и не выбирает либо не изменяет передаваемую информацию. \Cref{fig:conduit} сопоставляет эти условия с работой выходного ретранслятора Tor на уровне протокола ячеек.

\begin{figure}[htbp]
\centering
% SPDX-License-Identifier: GPL-3.0-or-later
% Copyright (C) 2026 Alessandro Greco
%% Русское альтернативное описание EPUB.
\altfig{Сравнение со статьёй 4 DSA не означает автоматического освобождения. Клиент использует RELAY\_BEGIN, узел exit открывает TCP, применяет правила и передаёт прикладные байты из RELAY\_DATA, а не сами ячейки Tor.}
% Требования mere conduit в сопоставлении с поведением exit
\begin{tikzpicture}[font=\sffamily\scriptsize,node distance=0.35cm,
  cond/.style={draw,rounded corners=2pt,text width=4.2cm,align=left,minimum height=0.8cm},
  tor/.style={draw=torviola,fill=torviola!8,rounded corners=2pt,text width=5.4cm,align=left,minimum height=0.8cm}]
  \node[cond] (a) {(a) не инициирует передачу};
  \node[cond,below=of a] (b) {(b) не выбирает её получателя};
  \node[cond,below=of b] (c) {(c) не выбирает и не изменяет информацию};
  \node[tor,right=1.2cm of a] (ta) {клиент посылает \texttt{RELAY\_BEGIN}; exit открывает TCP};
  \node[tor,right=1.2cm of b] (tb) {адрес задаёт клиент; policy может отказать};
  \node[tor,right=1.2cm of c] (tc) {байты из \texttt{RELAY\_DATA} в TCP-поток};
  \foreach \x in {a,b,c} \draw[<->,torgrigio] (\x) -- (t\x);
  \node[above=4pt of a,font=\sffamily\small] {Регламент (ЕС) 2022/2065, ст.~4};
  \node[above=4pt of ta,font=\sffamily\small] {Tor exit};
  \node[font=\sffamily\tiny,text=torgrigio,align=center,below=0.5cm of c,xshift=3.3cm] {Освобождение оценивается по обстоятельствам};
\end{tikzpicture}

\caption{Техническое сопоставление условий \emph{mere conduit} из ст.~4 Регламента (ЕС) 2022/2065 и работы выхода Tor, а не автоматическое юридическое заключение. Клиент указывает адрес в \texttt{RELAY\_BEGIN}, а выходной узел открывает TCP-соединение и может отклонить его согласно политике выхода. Прикладные байты из \texttt{RELAY\_DATA} передаются в поток TCP, а не в неизменённых ячейках Tor. Применимость освобождения зависит от всех правовых условий и обстоятельств.}
\label{fig:conduit}
\end{figure}

Важны ещё два положения. Статья~8 исключает \emph{общую} обязанность контролировать передаваемую информацию или активно искать признаки противоправной деятельности: оператор в общем случае не обязан без разбора проверять трафик; промежуточный ретранслятор также не может прочесть зашифрованное содержимое прикладного уровня в открытом виде. Однако статья~4(3) сохраняет полномочия судебных и административных органов обязывать поставщиков услуг предотвращать или прекращать нарушения.

\paragraph{Открытый вопрос.} Регламент применяется к \enquote{услугам информационного общества}, которые право ЕС определяет как услуги, \emph{обычно предоставляемые за вознаграждение}. Подпадает ли под это определение бесплатно работающий добровольческий ретранслятор? Существуют аргументы в пользу широкого толкования, а само освобождение отражает принцип, который суд мог бы учитывать: тот, кто лишь передаёт информацию, не должен отвечать за её содержимое. Однако правовое положение добровольческих ретрансляторов Tor в Италии не установлено здесь посредством прямо применимого судебного решения. Эта неопределённость отчасти объясняет, почему некоторые выходные ретрансляторы обслуживаются ассоциациями или организациями, уже предоставляющими сетевые услуги.

\subsection{Гражданская ответственность и уголовные расследования}

Освобождение касается ответственности за передаваемую информацию. Оно не препятствует тому, чтобы уголовное расследование началось с IP-адреса выходного ретранслятора. С точки зрения потерпевшего правонарушение, совершённое через Tor, выглядит исходящим от выходного ретранслятора. Следователи могут первоначально провести обыск или изъять компьютеры оператора. Юридический FAQ Electronic Frontier Foundation для операторов ретрансляторов, составленный по законодательству США, но полезный как запись практического опыта, сообщает, что организации не известны случаи осуждения оператора исключительно за эксплуатацию ретранслятора, однако описывает расследования и изъятия, вызванные неправильным пониманием Tor~\cite{efflegalfaq}. Снизить практический риск могут несколько мер:
\begin{itemize}
  \item Размещать выходной ретранслятор на отдельном сервере в дата-центре, а не дома или на личном рабочем компьютере.
  \item Явно указывать, что IP-адрес принадлежит выходному ретранслятору, например через информативную запись обратного DNS и информационную страницу, которую tor может отдавать через DirPort с помощью \opt{DirPortFrontPage}.
  \item Заранее уведомить провайдера связи и согласовать порядок обработки сообщений о злоупотреблениях.
  \item Объяснять работу Tor в ответах на жалобы и сохранять переписку.
\end{itemize}

\subsection{Хранение данных}

Статья~132 Кодекса Италии о защите персональных данных требует от подпадающих под норму \emph{провайдеров} хранить данные телематического трафика в течение двенадцати месяцев для целей расследования и преследования преступлений~\cite{codiceprivacy}; для телефонного трафика срок составляет двадцать четыре месяца. Пункт~5-\emph{bis} сохраняет специальное правило: статья~24 Закона №~167 от 20 ноября 2017 года увеличивает срок до \textbf{семидесяти двух месяцев} для расследований преступлений, указанных в статьях~51(3-\emph{quater}) и 407(2)(a) Уголовно-процессуального кодекса, включая терроризм и другие тяжкие преступления~\cite{legge167}. Соответствующая часть Кодекса о защите персональных данных применяется к услугам электронной связи, \emph{доступным населению через публичные сети связи} (статья~121).

Является ли добровольческий оператор ретранслятора таким провайдером — ещё один нерешённый вопрос, зависящий в том числе от способа предоставления сервиса. Операторам, особенно организациям, уже предоставляющим услуги связи, следует до запуска выходного ретранслятора получить правовую оценку. С технической точки зрения при рекомендуемых настройках журналирования tor обычно не хранит историю адресов назначения пользователей. Выходной ретранслятор видит адреса назначения соединений, которые он пересылает, но не знает исходный IP-адрес клиента. Фактически доступные данные зависят от роли и конфигурации ретранслятора: нельзя предполагать, что выходной ретранслятор способен устанавливать личности отдельных пользователей Tor.

\subsection{Журналы — это персональные данные}

IP-адреса относятся к персональным данным. Общий регламент по защите данных (GDPR) требует обрабатывать только необходимые данные и хранить их не дольше необходимого (статья~5, принцип минимизации данных). Для ретранслятора выбор, согласованный с безопасностью, — журналировать только на уровне \texttt{notice}, оставлять \opt{SafeLogging} включённым и избегать захвата сетевых пакетов (\cref{sec:traffico-log}).

% ---------------------------------------------------------------------------
\section{Исследователи}\label{sec:legale-ricerca}

В истории исследований Tor были этически сомнительные эксперименты, некоторые из которых обсуждаются в \cref{cap:27-casi-rete}. Исследователи запускали ретрансляторы, которые наблюдали или изменяли трафик реальных пользователей, чтобы продемонстрировать атаки. Помимо этических вопросов, подобные эксперименты могут затрагивать нормы уголовного права Италии, включая положения о незаконном перехвате компьютерных коммуникаций (статья~617-\emph{quater} Уголовного кодекса), а также правила защиты данных.

\textbf{Tor Research Safety Board} — группа исследователей, консультирующая по безопасности экспериментов и разработавшая соответствующие рекомендации~\cite{trsb}:
\begin{itemize}
  \item По возможности используйте тестовую сеть и атакуйте только собственный трафик.
  \item Собирайте только необходимые данные и только такие данные, которые было бы безопасно опубликовать.
  \item Защищайте собранные данные и экспериментальные инструменты.
  \item Уменьшайте точность данных округлением или добавлением шума.
  \item Учитывайте сторонние наборы данных, которые при объединении с вашими могут раскрыть дополнительную информацию.
  \item Задайтесь вопросом, ожидали ли пользователи, что соответствующие данные останутся приватными.
  \item Убедитесь, что ожидаемая польза оправдывает риски.
\end{itemize}
Совет просит исследователей объяснить, что они хотят узнать, как намерены это сделать, какие риски создаёт эксперимент и почему польза оправдывает эти риски. Эта книга следует той же дисциплине: лабораторные работы, создающие реальный трафик Tor, используют приватную тестовую сеть (\cref{sec:rp-perche}), а анализ общедоступных документов, например в \cref{cap:24-metrics}, проводится автономно по открытым агрегированным данным.

% ---------------------------------------------------------------------------
\section*{Итоги}
\begin{itemize}
  \item Само по себе использование Tor законно; деятельность через Tor регулируется обычным правом.
  \item Статья~4 DSA освобождает подпадающие под неё сервисы, которые лишь передают чужую информацию, от ответственности при соблюдении определённых условий. Принципы работы выходного ретранслятора важны для оценки этих условий, но вопрос о том, подпадает ли добровольческий ретранслятор под Регламент, остаётся юридически открытым.
  \item Освобождение не исключает расследований и изъятия оборудования. Выделенные серверы, явное обозначение выходного ретранслятора и договорённости с провайдерами связи могут уменьшить практические проблемы.
  \item Требования к хранению данных применяются к подпадающим под норму общедоступным провайдерам связи; применимость к добровольческим ретрансляторам требует правовой оценки конкретного случая.
  \item Исследования Tor по возможности должны избегать экспериментов с реальными пользователями: используйте тестовые сети, минимизируйте данные, публикуйте только безопасные агрегаты и сопоставляйте пользу с рисками.
\end{itemize}

% ---------------------------------------------------------------------------
\section*{Упражнения}

\begin{esercizio}[1]\label{es:25-condizioni}
Для каждого из трёх условий статьи~4 DSA опишите изменение в поведении выходного ретранслятора, которое нарушило бы соответствующее условие.
\end{esercizio}

\begin{esercizio}[1]\label{es:25-strumento}
На примере объясните различие между законностью инструмента и законностью его использования, применив его и к пользователю Tor, и к оператору ретранслятора.
\end{esercizio}

\begin{esercizio}[2]\label{es:25-segnalazione}
Составьте ответ на сообщение о злоупотреблениях, полученное от провайдера связи выходного ретранслятора после жалобы сайта на повторяющиеся попытки доступа с вашего IP-адреса. Что можно объяснить и какие практические меры предпринять?
\end{esercizio}

\begin{esercizio}[2]\label{es:25-esperimento}
Исследователь предлагает подсчитывать доменные имена, наблюдаемые на выходном ретрансляторе, чтобы оценить, сколько пользователей Tor посещают итальянские новостные сайты. Оцените предложение по рекомендациям Tor Research Safety Board и предложите более безопасную альтернативу.
\end{esercizio}

\begin{esercizio}[3]\label{es:25-ricerca}
Найдите судебные решения, касающиеся операторов ретрансляторов Tor, в судах Италии и других европейских стран. Для каждого кратко изложите юридический вопрос и результат, ссылаясь на первичные источники.
\end{esercizio}

% ---------------------------------------------------------------------------
\section*{Дополнительная литература}
\begin{itemize}
  \item Регламент (ЕС) 2022/2065, в особенности статьи~4, 8 и~9~\cite{dsa2022}.
  \item Законодательный декрет 196/2003, статьи~121 и~132, в сводном тексте, опубликованном Итальянским органом по защите персональных данных~\cite{codiceprivacy}.
  \item Electronic Frontier Foundation, \titolo{Legal FAQ for Tor Relay Operators}~\cite{efflegalfaq}.
  \item Сайт Tor Research Safety Board~\cite{trsb}.
\end{itemize}
